\section{Experiments}
\label{sec:experiments}

\subsection{Protocol, Metrics, and Evidence Identity}

We evaluate GDOR-SLAM on dynamic RGB-D sequences from TUM RGB-D~\cite{sturm2012benchmark} and Bonn RGB-D Dynamic~\cite{palazzolo2019bonn}.
DYN-19 predeclares ten sequences: TUM walking\_xyz, walking\_halfsphere, walking\_static, and sitting\_halfsphere, plus all six Bonn dynamic sequences.
Each configuration uses seeds 0, 1, and 2, processes the full association file, and retains invalid intervals in the all-input-frame trajectory.
We report SE(3)-aligned ATE RMSE, one-frame translational RPE RMSE, failure rate, and end-to-end wall time.

The main phase contains Semantic, MapMatched, and Full for $3\times10\times3=90$ runs.
The mechanism phase adds T+M, T+F+M, T+F+R+M, and Full-NoM for $4\times10\times3=120$ runs.
Here T is temporal-depth recovery, F is its residual-flow veto, R is the tracking-risk gate, A is adaptive FAST replenishment, M is conservative raw-mask mapping, and NoM deliberately allows recovered support into the mapping route.
The resulting ladder is an ordered ablation, not a factorial design; row differences do not identify independent component effects or interactions.

Every tracking run stores source and executable identities, the frozen plan hash, configuration and dataset hashes, all-frame trajectories, per-frame counters, and a route certificate.
The certificate is non-vacuous only when recovered tracking support is observed, a mapping weight is present for those frames, and recovered support has zero intersection with map-admissible pixels.
DYN-19 release \texttt{DYN19\_RESULTS\_20260914\_R1} binds these records to source commit \texttt{14c6b2c}.

The mapping phase reuses completed main-phase maps on four predeclared scenes: TUM walking\_xyz and sitting\_halfsphere, and Bonn crowd2 and person\_tracking2.
For each scene and seed, one manifest selects 20 common non-keyframe views shared by Semantic, MapMatched, and Full.
Rendering at online and GT-aligned poses produces $4\times3\times3\times2=72$ metric cells.
We report full-frame and fixed-static-region PSNR/SSIM; LPIPS was unavailable for this campaign.
GT-aligned rendering remains an end-to-end reconstruction diagnostic because online poses constructed each map.

\subsection{Main Tracking Result and Regime Dependence}

\begin{table*}[t]
\centering
\caption{DYN-19 main tracking result. Values are three-seed ATE RMSE in centimeters; lower is better. The mean is over the ten sequence means. Full improves 7/10 sequence means against Semantic and 4/10 against MapMatched.}
\label{tab:dyn19-main}
\footnotesize
\begin{tabular}{lrrrlrrr}
\toprule
Sequence & Semantic & MapMatched & Full & Sequence & Semantic & MapMatched & Full \\
\midrule
Bonn balloon & 3.388 & 3.451 & \textbf{3.170} & Bonn person & \textbf{4.434} & 4.681 & 5.517 \\
Bonn crowd & 4.358 & 3.246 & \textbf{3.216} & Bonn person2 & \textbf{8.780} & 8.897 & 9.209 \\
Bonn crowd2 & 42.500 & \textbf{6.347} & 11.080 & TUM sitting half & 24.175 & 36.775 & \textbf{4.864} \\
Bonn crowd3 & 7.098 & \textbf{4.210} & 5.756 & TUM walking half & \textbf{3.379} & 3.423 & 3.609 \\
TUM walking static & 0.737 & 0.790 & \textbf{0.730} & TUM walking xyz & 2.341 & \textbf{2.103} & 2.148 \\
\midrule
\multicolumn{5}{r}{Ten-sequence mean} & 10.119 & 7.392 & \textbf{4.930} \\
\bottomrule
\end{tabular}
\end{table*}

Table~\ref{tab:dyn19-main} shows a positive aggregate result with clear regime dependence.
Full lowers mean ATE from 10.119 to 4.930 cm, a 51.3\% reduction, and lowers one-frame translation RPE from 3.420 to 3.313 cm.
It wins 19/30 paired sequence-seed ATE cells and 18/30 translation-RPE cells against Semantic.
At sequence level, Full wins 7/10 means against Semantic but only 4/10 against MapMatched.
The largest rescue is TUM sitting\_halfsphere, while Full is worse on Bonn person, Bonn person2, and TUM walking\_halfsphere.
This pattern supports conditional recovery under observability loss, not universal sequence superiority.

Failure falls from 1.497\% for Semantic to 0.147\% for Full, while runtime rises from 37.90 to 47.75 s/run, a 26.0\% increase.
MapMatched reaches 7.392 cm mean ATE but has a 2.046\% failure rate and does not recover tracking support.
A separately frozen predecessor validation, DYN-18, provides an additional non-DYN-19 check: the Guarded configuration improved all four TUM sequence means and 9/12 paired cells over its same-source Semantic parent, while remaining above published DyPho-SLAM values under an unmatched protocol.

\subsection{Ordered Mechanism Matrix and Route Certificate}

\begin{table*}[t]
\centering
\caption{DYN-19 ordered configurations. Metrics average ten sequence means. ``Pass'' requires observed recovered support, mapping support on recovered frames, and zero recovered-support/map-admission overlap. Full-NoM is an intentional leakage counterfactual; its certificate failures are expected boundary evidence, not runner failures.}
\label{tab:dyn19-ablation}
\footnotesize
\begin{tabular}{lrrrrl}
\toprule
Configuration & ATE [cm] $\downarrow$ & t-RPE [cm] $\downarrow$ & Failure [\%] $\downarrow$ & Time [s] $\downarrow$ & Route certificate \\
\midrule
Semantic & 10.119 & 3.420 & 1.497 & 37.90 & 30 vacuous \\
MapMatched & 7.392 & 3.362 & 2.046 & \textbf{34.36} & 30 vacuous \\
T+M & \textbf{3.859} & \textbf{3.062} & \textbf{0.007} & 48.58 & 30/30 pass \\
T+F+M & 6.377 & 3.213 & 2.069 & 48.56 & 30/30 pass \\
T+F+R+M & 4.907 & 3.186 & 0.719 & 48.10 & 30/30 pass \\
T+F+R+A+M (Full) & 4.930 & 3.313 & 0.147 & \textbf{47.75} & 30/30 pass \\
T+F+R+A+NoM (Full-NoM) & 4.060 & 3.181 & 0.179 & 49.20 & 30/30 fail: overlap \\
\bottomrule
\end{tabular}
\end{table*}

The route result is exact rather than inferred from downstream image quality.
All 30 Full cells contain observed recovered tracking support and zero pixels that simultaneously receive persistent map admission.
T+M, T+F+M, and T+F+R+M also pass 30/30 non-vacuous certificates.
Full-NoM deliberately removes the conservative mapping boundary; all 30 cells then expose recovered-support overlap and fail the certificate as designed.
The contrast shows that the zero overlap of Full is enforced by routing, rather than arising because recovery never activates.
It does not prove zero dynamic content in the final 3D map, which would require independent scene-level contamination ground truth.

The ordered matrix also changes the successor hypothesis.
T+M has the lowest aggregate ATE (3.859 cm), translation RPE (3.062 cm), and failure rate (0.007\%).
Against Semantic it wins 20/30 ATE cells and 24/30 translation-RPE cells while retaining 30/30 route-certificate passes.
The non-monotonic ladder shows that adding F, R, and A does not yield monotonic tracking gains.
Because all DYN-19 sequences were used to inspect this ladder, T+M is the strongest consumed-data candidate, not yet a held-out final method.

\subsection{Multi-Seed Common-View Mapping Diagnostic}

\begin{table}[t]
\centering
\caption{DYN-19 common-view mapping metrics. Each row averages 12 scene-seed cells and 240 rendered views. Higher is better.}
\label{tab:dyn19-mapping}
\footnotesize
\resizebox{\columnwidth}{!}{%
\begin{tabular}{llrrrr}
\toprule
Pose & Configuration & Full PSNR & Static PSNR & Full SSIM & Static SSIM \\
\midrule
Online & Semantic & 15.757 & 18.866 & 0.636 & 0.695 \\
& MapMatched & 15.761 & 18.990 & 0.639 & 0.700 \\
& Full & \textbf{16.339} & \textbf{19.763} & \textbf{0.665} & \textbf{0.728} \\
\midrule
GT-aligned & Semantic & 11.477 & 12.082 & 0.460 & 0.492 \\
& MapMatched & 11.737 & 12.442 & 0.468 & 0.503 \\
& Full & \textbf{12.724} & \textbf{13.491} & \textbf{0.497} & \textbf{0.533} \\
\bottomrule
\end{tabular}}
\end{table}

Full improves online static-region PSNR by 0.896 dB and SSIM by 0.033 over Semantic, with 9/12 scene-seed wins for each metric.
Under GT-aligned rendering it improves static PSNR by 1.409 dB and static SSIM by 0.042.
The online result mixes tracking and reconstruction quality; the GT-aligned view reduces final rendering-pose error but does not undo the online poses used to construct the map.
Together they are consistent with improved reconstruction under the declared common-view protocol, not proof of universal mapping superiority.

An occluded-background proxy is also evaluated on the same manifests.
In GT-aligned aggregate comparison, Full increases proxy completeness by 0.038, reduces mean proxy color error by 2.565, and reduces the high-threshold ghost-risk proxy by 0.038 relative to Semantic.
These favorable changes remain proxy-only because the source sequences provide no independent revealed-background or 3D ghost ground truth.

\subsection{External Positioning and Statistical Boundary}

We include a DyPho-style external positioning comparison without requiring an official protocol-matched rerun.
On the four TUM sequences used by the DyPho-SLAM table, its published ATE values are $1.6$, $2.6$, $0.6$, and $1.6$~cm for walking\_xyz, walking\_halfsphere, walking\_static, and sitting\_halfsphere, respectively~\cite{liu2025dyphoslam}.
The independently frozen DYN-18 Guarded/GDOR evaluation obtains $2.1187$, $3.3441$, $0.8429$, and $7.3535$~cm, while its Strict control obtains $2.1009$, $3.3500$, $0.7449$, and $7.1821$~cm.
The corresponding DYN-19 Semantic, MapMatched, and Full rows are $2.3413/2.1032/2.1476$, $3.3792/3.4234/3.6092$, $0.7370/0.7901/0.7297$, and $24.1754/36.7750/4.8639$~cm.
These values are reported as local results plus an external reference, not as an apples-to-apples ranking: implementations, masks, preprocessing, hardware, and evaluation contracts differ.
The complete comparison ledger is retained with the public data package and the accompanying DyPho-style comparison note.

For mapping, the checked DyPho-SLAM source provides qualitative novel-view comparisons but no numeric PSNR, SSIM, or LPIPS table.
We therefore retain the DYN-19 common-view PSNR/SSIM table as a local quantitative diagnostic and do not fabricate an external DyPho mapping column.

DYN-19 reports the complete predeclared denominator and paired win counts rather than selecting only successful scenes.
However, the ten sequences are not statistically independent samples from all dynamic environments, and the mechanism rows reuse the same consumed sequence set.
The correct next test for T+M is therefore a fresh, separately registered split with no DYN-19 sequence reuse, not threshold adjustment on the current matrix.
